\section{世界模型：从公式、控制到机器人大脑}

本期最硬的技术部分，是谢赛宁对世界模型的定义。他没有把 world model 当成流行词，而是回到控制、model-based reinforcement learning 和认知科学里的基本问题：给定当前状态和动作，预测下一状态；再用这种预测指导 planning 和 decision making。

\[
s_{t+1} = f(s_t, a_t)
\]

其中，$s_t$ 表示当前时刻系统或环境的状态；$a_t$ 表示智能体采取的 action 或 intervention；$f$ 是学到的 transition / predictive function；$s_{t+1}$ 是执行动作后可能进入的下一状态。真正困难的部分不在公式，而在：怎样从高维感知中得到合适的 $s_t$，怎样让 $f$ 学到物理规律，怎样把预测用于行动而不是只生成漂亮样本。

\begin{figure}[H]
\centering
\includegraphics[width=0.96\textwidth]{world-model-stack.png}
\caption{世界模型的最小闭环：感知、状态表征、转移预测、规划决策、真实反馈和学习更新。}
\end{figure}

\begin{knowledgebox}{读图：世界模型不是一个单独模块}
图中最重要的是闭环。感知输入先被压缩成任务相关 state；模型根据 state 和 action 预测下一状态；规划器比较不同未来轨迹；动作执行后得到真实反馈；误差和失败再回流训练。少掉任何一环，world model 都会退化成静态生成器、问答器或离线表征模型。
\end{knowledgebox}

\subsection{Model Predictive Control 的直觉}

前面公式说明了 world model 的最小形式；现在需要回答它怎样变成行动。MPC 提供了一个最朴素也最有教学价值的答案：用模型想象未来，再只执行当前最该做的一步。

访谈中提到 model predictive control。它的基本做法是：在当前时刻，用模型 roll out 多个未来 action sequence，计算每条轨迹的 cost，选择 cost 最低的序列，执行第一步，然后下一时刻重新规划。这不是 LLM 式的“说出理由”，而是以预测未来状态为基础的行动选择。

\begin{knowledgebox}{术语消化：世界模型相关概念}
\begin{tabularx}{\textwidth}{p{0.23\textwidth}p{0.30\textwidth}X}
\toprule
术语 & 一句话解释 & 在本期中的作用 \\
\midrule
State & 描述系统当前状态的最小充分信息 & 连接表征学习与行动预测。 \\
Action / Intervention & 智能体对环境施加的动作或干预 & 让预测从观察变成决策。 \\
Transition Function & 从当前 state 和 action 预测下一 state 的函数 & 世界模型的核心形式。 \\
Planning & 比较未来轨迹并选择动作 & 让模型预测服务目标。 \\
MPC & Model Predictive Control，滚动预测与控制 & 展示世界模型如何用于控制。 \\
Model-based RL & 用环境模型辅助强化学习 & 说明 world model 与 RL 的历史关系。 \\
\bottomrule
\end{tabularx}
\end{knowledgebox}

\subsection{World model 是目的，不是单一算法}

谢赛宁特别强调，世界模型不好定义，是因为它不是一个算法名，而是一个目的。语言模型、video diffusion model、3D representation、robotics、VLA、model-based RL，都可能从不同方向走向 world model。争论“谁才是真的 world model”在短期内有意义，但长期看更重要的问题是：这个系统是否能理解物理世界、保留相关 memory、reason and plan、做 counterfactual / causal inference，并且 controllable and safe。

\begin{importantbox}{世界模型的五项能力}
面向物理世界的 world model 至少要具备：physical world understanding、large associated memory、reasoning / planning、counterfactual or causal inference、controllability and safety。它不是“视频生成得像不像”这么单一的指标。
\end{importantbox}

\subsection{本章小结}

世界模型的底层定义很简单：预测 action 后的 state。难点在于 state 的表征、预测的物理接地、planning 的可用性和反馈闭环。把 world model 当成目标，而不是某个模型架构，可以避免被短期名词战带偏。

